\begin{claim}{D.41}{Proved}{The global multiplier of a negative transposition}
An involution of the full three-cusp loop-outer group whose
puncture permutation is a transposition and whose $2$-adic
multiplier is $-1$ has the global multiplier
$-1\in\widehat\Z^\times$.
Together with D.34 and D.36, this identifies the multiplier
of every involution in that full group with its global
$2$-adic sign.  Consequently every finite subgroup generated
by involutions has the usual global orientation character.
\end{claim}
\begin{proof}
Label the original peripheral generators as
$x,y,z=(xy)^{-1}$ so that the involution fixes $x$ and
exchanges $y,z$.  Write its common multiplier as $\lambda$.
It satisfies $\lambda^2=1$, so every coordinate $\lambda_p$
is $1$ or $-1$.  The peripheral-normalizer construction
(D.37d)--(D.37e) supplies an actual order-two automorphism
$\alpha$ of $F=\widehat F(x,y)$ representing it and satisfying
$\alpha(x)=x^\lambda$.

Suppose a positive odd-prime coordinate exists, and let
$q$ be the least odd prime with $\lambda_q=1$.  Use the
exact cyclic quotient
\[
 \chi:F\longrightarrow C_q,\qquad
       \chi(x)=-2,\quad\chi(y)=\chi(z)=1,\qquad N=\ker\chi .
                                                        \tag{D.41a}
\]
The relation $xyz=1$ is respected.  Since $\lambda=1$
modulo $q$ and $\alpha$ exchanges the last two peripheral
classes, $\chi\alpha=\chi$.  The cover corresponding to
$N$ is the completion of the specified connected cyclic
cover of the pair of pants.  Each of its three original
cusps has one lift of degree $q$: all three displayed
monodromies are nonzero in $C_q$.  Its Euler characteristic
is $-q$, so its genus is
\[
                             g=(q-1)/2.
                                                        \tag{D.41b}
\]

Fill the two cusp lifts represented by $y^q,z^q$.
The quotient and its deck extension are explicit:
\[
\begin{split}
 R&=F/\overline{\langle\!\langle y^q,z^q\rangle\!\rangle}_F
       =\langle y,z\mid y^q=z^q=1\rangle_{\rm prof}
       =C_q\amalg C_q,\\
 K&=\ker(R\longrightarrow C_q,\ y,z\longmapsto1)\\
  &=N/\overline{\langle\!\langle y^q,z^q
                                      \rangle\!\rangle}_N
       \cong\widehat\pi_{g,1}.
\end{split}                                             \tag{D.41c}
\]
Here $x=z^{-1}y^{-1}$, and the one remaining cusp is
$c=x^q$.  To verify the equality of normal subgroups
used in (D.41c), note $F=N\langle y\rangle=N\langle z\rangle$.
Conjugating $y^q$ or $z^q$ by any element of $F$ therefore
gives an $N$-conjugate of the same element.  Thus its
closed $F$-normal closure is its closed $N$-normal closure.
The completion comparison is also exact: start with the
discrete cyclic surface cover, fill those two boundaries,
and complete the displayed quotient presentation.
Finite-index subgroup topology identifies its completed
kernel with the specified kernel in $C_q\amalg C_q$.
The inverse presentation substitutions are the original
$y,z$ and $x=z^{-1}y^{-1}$.

The automorphism $\alpha$ preserves this normal subgroup,
because it exchanges the two filled peripheral classes.
It therefore induces an actual automorphism of $R$,
preserves $K$, and induces an involution $\beta$ of $K$
with
\[
                           \beta(c)=c^\lambda.
                                                        \tag{D.41d}
\]
Let $T$ be conjugation by $y$ on $K$.  In $R$, $y^q=1$,
so $T^q=1$ as an actual automorphism.  The equality
$\chi\alpha=\chi$ gives $\alpha(y)y^{-1}\in K$; hence
$\beta T\beta^{-1}$ and $T$ differ by an inner
automorphism of $K$.  In particular their maps on
abelianization commute.

We compute that abelianization as a deck module, without
requiring a comparison of traces.  Reidemeister--Schreier
rewriting with representatives $1,y,\ldots,y^{q-1}$ gives
\[
 u_i=y^i z y^{-(i+1)}\quad(i\in\Z/q),\qquad
 u_0u_1\cdots u_{q-1}=1,\qquad T(u_i)=u_{i+1}.
                                                        \tag{D.41e}
\]
Exponents of $y$ here are interpreted in $R$, where
$y^q=1$.  The first $q-1$ elements form a free profinite
basis; eliminate the last by
$u_{q-1}=(u_0\cdots u_{q-2})^{-1}$.
These substitutions give both maps in the Schreier
presentation.  The remaining boundary also has the exact
word
\[
 c=x^q=\prod_{j=0}^{q-1}u_{\,q-1-2j}^{-1},
                       \qquad\text{indices modulo }q.
                                                        \tag{D.41f}
\]
Indeed $z=u_0y$ gives
$x=z^{-1}y^{-1}=y^{-1}u_0^{-1}y^{-1}
 =u_{q-1}^{-1}y^{-2}$.
Multiplying its $q$ successive translates by $y^{-2}$
gives exactly (D.41f), because $y^{2q}=1$.

For any prime $p$ put $V=H_1(K,\Q_p)$ and let $A$ be
the map induced by $\beta$.  From (D.41e) we obtain
the exact module isomorphism
\[
 V\cong
 \Q_p[C_q]\big/\Q_p(1+T+\cdots+T^{q-1})
 \cong\Q_p[X]/\Phi_q(X),\qquad
              A^2=1,\quad AT=TA.
                                                        \tag{D.41g}
\]
Every irreducible cyclotomic factor occurs with
multiplicity one.

The one-boundary surface structure in (D.41c) supplies
a nondegenerate alternating pairing, and its relation
to (D.41d) is algebraic.  In a surface free basis the
boundary word is the product of the $g$ commutators,
up to inversion and conjugation.  Its degree-two image
in the class-two free pro-$p$ quotient is a nondegenerate
bivector $\omega\in\bigwedge^2 V$.
Equation (D.41d) gives
$(\bigwedge^2 A)\omega=\lambda_p\omega$.
The deck map preserves its peripheral exponent, so
$(\bigwedge^2 T)\omega=\omega$.
If $\Omega$ is the nonsingular skew matrix of $\omega$,
these are precisely
\[
 A\Omega A^{\mathsf t}=\lambda_p\Omega,\qquad
 T\Omega T^{\mathsf t}=\Omega.
\]
Consequently the alternating form with matrix
$\Omega^{-1}$ satisfies
\[
 \langle Av,Aw\rangle=\lambda_p\langle v,w\rangle,
 \qquad
 \langle Tv,Tw\rangle=\langle v,w\rangle.
                                                        \tag{D.41h}
\]
This proves the multiplier comparison also at $p=2$:
one may first take the integral commutator calculation
in each finite class-two quotient and then tensor with
$\Q_2$.  It uses no geometric realization of $\beta$.

Now let $p\ne q$ be a prime with
$d=\operatorname{ord}_q(p)$ even.  The irreducible
factors of $\Phi_q$ over $\Q_p$ have root sets
the Frobenius orbits
\[
                  \{\zeta^{j p^k}:0\leq k<d\},
                     \qquad j\in(\Z/q)^\times.
                                                        \tag{D.41i}
\]
For completeness these roots lie in the unramified
extension with residue field $\F_{p^d}$: the primitive
$q$th roots in that field lift uniquely since $p\ne q$,
and Frobenius lifts their map $z\mapsto z^p$.
The least such degree is $d$, proving the factor
description.  Since $d$ is even, $p^{d/2}=-1$ modulo
$q$, and each of these root sets is closed under
inversion.

Write $V=\bigoplus V_i$ for those irreducible summands.
By (D.41g), the commutant of $T$ is the product of their
fields $\Q_p(\zeta^j)$.  Therefore $A|_{V_i}$ is
multiplication by an element of a field whose square
is one, and so is the scalar $1$ or $-1$ on that entire
summand.  Each $V_i$ is nondegenerate for the form in
(D.41h).  To see this explicitly after a splitting-field
extension, a $T$-eigenvector with eigenvalue $\zeta^j$
can pair only with the eigenvalue $\zeta^{-j}$, since
the pairing is $T$-invariant.  Both belong to $V_i$ by
(D.41i); different summands are orthogonal.  Global
nondegeneracy thus restricts to nondegeneracy on $V_i$.
Multiplication by either scalar $1$ or $-1$ preserves
that restricted form.  Comparing with (D.41h) gives
\[
 \lambda_p=1
       \quad\text{whenever }p\ne q
                    \text{ and }\operatorname{ord}_q(p)
                    \text{ is even}.
                                                        \tag{D.41j}
\]

There is a prime divisor $p$ of $q-1$ having this even
order.  Otherwise every prime factor of $q-1$ would
lie in the odd-order subgroup of the cyclic group
$(\Z/q)^\times$.  Their product, with multiplicities,
would also have odd order.  But that product is
$q-1=-1$ modulo $q$, of order two.  This contradiction
proves the assertion.  Such a prime satisfies $p<q$.
For $p=2$ the hypothesis gives $\lambda_p=-1$; for an
odd $p<q$ the minimal choice of $q$ gives the same
value.  Both contradict (D.41j).

Thus no positive odd-prime coordinate exists, and
$\lambda=-1$ globally.  D.34 handles a pure
involution, while D.36 handles an involution with
positive $2$-adic multiplier.  All involutions are
therefore accounted for.  Multiplicativity proves
the last assertion for every subgroup generated by
them.  This conclusion determines these actual
multiplier homomorphisms; it does not assert geometric
conjugacy of the finite actions or settle the
multiplier of an isolated order-three subgroup.

The program \texttt{code/part\_c/negative\_multiplier\_d41.py}
checks the exact Schreier boundary word, the induced
cyclotomic deck matrices and the nondegenerate
degree-two boundary form in examples, and checks
the prime-factor and Frobenius-orbit calculation.
The displayed argument proves the assertion for
every odd prime, rather than only for the executed
range.
\end{proof}
